\documentclass[11pt]{article}
\usepackage{amsmath, amsthm, amssymb}
\usepackage[margin=1.1in]{geometry}
\usepackage{hyperref}

\newtheorem{theorem}{Theorem}
\newtheorem{proposition}[theorem]{Proposition}
\newtheorem{corollary}[theorem]{Corollary}
\newtheorem{lemma}[theorem]{Lemma}
\theoremstyle{definition}
\newtheorem*{conjecture}{Conjecture}
\theoremstyle{remark}
\newtheorem*{remark}{Remark}

\newcommand{\Frob}{\operatorname{Frob}}
\newcommand{\Ind}{\operatorname{Ind}}
\newcommand{\Res}{\operatorname{Res}}
\newcommand{\Hom}{\operatorname{Hom}}
\newcommand{\sgn}{\operatorname{sgn}}
\newcommand{\ch}{\operatorname{ch}}
\newcommand{\triv}{\mathrm{triv}}
\newcommand{\fdeg}{\widetilde{f}}
\newcommand{\qbinom}[2]{\left[{#1 \atop #2}\right]_{\!t}}

\title{Period-3 divisibility of the $S_2$-isotypic\\multiplicities of $R_{(k,k)}$}
\author{Clio}
\date{2026-07-31}

\begin{document}
\maketitle

\begin{abstract}
Let $R_{(k,k)} = R_{2k}^{S_k \times S_k}$ denote the parabolic coinvariant
algebra of $S_{2k}$ at the composition $\alpha = (k,k)$, viewed as a
graded $S_2$-module via the block-swap action. Let
$m_{(2)}(t), m_{(1,1)}(t) \in \mathbb{Z}[t]$ be the trivial and sign
isotypic Hilbert polynomials. We prove
\[
[3]_t \mid m_{(2)}(t) \quad\text{and}\quad [3]_t \mid m_{(1,1)}(t)
\qquad \iff \qquad k \equiv 2 \pmod 3.
\]
The proof reduces both $m_\pi(t)$ to alternating partial sums of
$q$-binomial coefficients via the Kirillov--Reshetikhin two-row fake
degree, evaluates at $\zeta_3$ via the ``correct'' $q$-analogue of Lucas'
theorem, and observes that at $k \equiv 2 \pmod 3$ every full block of
three consecutive $i$'s contributes zero, whereas at $k \equiv 0, 1
\pmod 3$ a nonzero residue survives.
\end{abstract}

\section{Setup and statement}

Fix $k \ge 1$ and write $n = 2k$, $\alpha = (k, k)$. Let
$R_n = \mathbb{C}[x_1, \ldots, x_n] / \langle e_1, \ldots, e_n\rangle$
be the coinvariant algebra of $S_n$, and
$R_\alpha := R_n^{S_\alpha}$ the $S_\alpha = S_k \times S_k$
invariants. The normaliser $N_{S_n}(S_\alpha) = S_k \wr S_2$ acts on
$R_\alpha$; the residual $S_2 = N/S_\alpha$ acts by block-swap.

\begin{theorem}[\cite{Wreath-r=2}, Theorem D at $r=2$]\label{thm:D}
The graded $S_2$-isotypic Hilbert polynomials of $R_\alpha$ are
\begin{align*}
m_{(2)}(t) &= \sum_{i=0}^{\lfloor k/2\rfloor} \fdeg_{(2k-2i,\,2i)}(t), \\
m_{(1,1)}(t) &= \sum_{i=0}^{\lfloor (k-1)/2\rfloor} \fdeg_{(2k-2i-1,\,2i+1)}(t),
\end{align*}
where $\fdeg_\lambda(t) = t^{n(\lambda)} [n]_t! / \prod_{c \in \lambda} [h(c)]_t$
is the fake degree of $\lambda$.
\end{theorem}

Our result:

\begin{theorem}[Period-$3$ divisibility, $r=2$]\label{thm:main}
For every $k \ge 1$,
\[
[3]_t \mid m_{(2)}(t) \text{ and } [3]_t \mid m_{(1,1)}(t)
\quad \iff \quad k \equiv 2 \pmod 3.
\]
\end{theorem}

\section{The two-row fake degree}

\begin{lemma}\label{lem:two-row}
For $a \ge b \ge 0$ integers, letting $\qbinom{n}{k} = [n]_t!/([k]_t! [n-k]_t!)$
denote the Gaussian binomial (with $\qbinom{n}{-1} := 0$),
\[
\fdeg_{(a,b)}(t) \;=\; \qbinom{a+b}{b} \;-\; \qbinom{a+b}{b-1}.
\]
\end{lemma}

\begin{proof}
For $\lambda = (a, b)$, $n(\lambda) = b$, and the hook lengths are
$\{1, 2, \ldots, a+1\} \setminus \{a - b + 1\}$ in the first row and
$\{1, 2, \ldots, b\}$ in the second row. Thus
\[
\fdeg_{(a,b)}(t) \;=\; t^b \cdot \frac{[a+b]_t!}{[b]_t! \cdot [a+1]_t! / [a-b+1]_t}
\;=\; \frac{t^b \cdot [a-b+1]_t \cdot [a+b]_t!}{[b]_t! \cdot [a+1]_t!}.
\]
On the other hand, from $\qbinom{a+b}{b} = [a+b]_t!/([b]_t! [a]_t!)$ and
$\qbinom{a+b}{b-1} = [a+b]_t!/([b-1]_t! [a+1]_t!)$,
\begin{align*}
\qbinom{a+b}{b} - \qbinom{a+b}{b-1}
&= \frac{[a+b]_t!}{[b]_t![a+1]_t!}\bigl([a+1]_t - [b]_t\bigr).
\end{align*}
The bracket telescopes: $[a+1]_t - [b]_t = t^b + t^{b+1} + \cdots + t^a = t^b [a-b+1]_t$.
Substituting matches the fake-degree expression above.
\end{proof}

\section{The telescoping identity}

Set $c_j := \qbinom{2k}{j}$ (so $c_j \in \mathbb{Z}[t]$) and $c_{-1} := 0$. Define the alternating partial sum
\[
S_j(t) \;:=\; \sum_{i=0}^{j} (-1)^i c_i \;=\; \sum_{i=0}^{j} (-1)^i \qbinom{2k}{i}.
\]

\begin{proposition}\label{prop:telescope}
For every $k \ge 1$,
\[
m_{(2)}(t) \;=\; S_{j}(t), \qquad m_{(1,1)}(t) \;=\; -\,S_{j'}(t),
\]
where
\[
(j, j') \;=\; \begin{cases}(k,\ k-1) & \text{$k$ even},\\
(k-1,\ k) & \text{$k$ odd}.\end{cases}
\]
Consequently
\[
m_{(2)}(t) + m_{(1,1)}(t) \;=\; \qbinom{2k}{k}.
\]
\end{proposition}

\begin{proof}
By Lemma~\ref{lem:two-row}, $\fdeg_{(2k-2i,\,2i)}(t) = c_{2i} - c_{2i-1}$, so
\[
m_{(2)}(t) = \sum_{i=0}^{K}(c_{2i} - c_{2i-1})
= c_0 - c_1 + c_2 - c_3 + \cdots + (-1)^{2K} c_{2K}
= \sum_{i=0}^{2K}(-1)^i c_i
= S_{2K}(t),
\]
with $K = \lfloor k/2 \rfloor$; and $2K = k$ if $k$ even, $2K = k-1$ if $k$ odd.

Similarly $\fdeg_{(2k-2i-1,\,2i+1)}(t) = c_{2i+1} - c_{2i}$, so
\[
m_{(1,1)}(t) = \sum_{i=0}^{K'} (c_{2i+1} - c_{2i})
= -\sum_{i=0}^{2K'+1}(-1)^i c_i
= -S_{2K'+1}(t),
\]
with $K' = \lfloor (k-1)/2 \rfloor$; and $2K' + 1 = k$ if $k$ odd, $k-1$ if $k$ even.

For the final identity, note that $S_k(t) - S_{k-1}(t) = (-1)^k c_k$.
If $k$ is even, $m_{(2)} + m_{(1,1)} = S_k - S_{k-1} = c_k$; if $k$ is odd,
$m_{(2)} + m_{(1,1)} = S_{k-1} - S_k = -(-1)^k c_k = c_k$. Either way the
sum equals $c_k = \qbinom{2k}{k}$. (Equivalently: this is
$h_2[h_k] + e_2[h_k] = h_k^2$ on the level of fake degrees, using the
Pieri expansion $h_k^2 = \sum_{j=0}^k s_{(k+j, k-j)}$ and the MacMahon
identity $\sum_j \fdeg_{(k+j, k-j)}(t) = \binom{2k}{k}_t$.)
\end{proof}

\begin{remark}
The identity $\sum_{j=0}^{k} \fdeg_{(k+j, k-j)}(t) = \binom{2k}{k}_t$
is the classical MacMahon formula for the Hilbert series of the parabolic
coinvariant $R_{(k, k)} = R_{2k}^{S_k \times S_k}$; it is also
Theorem~D's total Hilbert series identity at $r = 2$.
\end{remark}

\section{The $q$-Lucas evaluation}

Let $\omega = e^{2\pi i/3}$ so $[3]_t = 1 + t + t^2 = \Phi_3(t)$ has roots
$\{\omega, \bar\omega\}$. Since $[3]_t = \Phi_3(t)$ is irreducible over
$\mathbb{Q}[t]$, for $f \in \mathbb{Q}[t]$:
\[
[3]_t \mid f(t) \iff f(\omega) = 0.
\]

We use the following ``correct'' form of the $q$-Lucas theorem.

\begin{lemma}[$q$-Lucas at $\zeta_d$, Désarménien 1982]\label{lem:qlucas}
For any positive integer $d$, any primitive $d$-th root of unity $\zeta$,
and any nonnegative integers $n \ge k$,
\[
\qbinom{n}{k}\!\bigg|_{t = \zeta} \;=\;
\binom{\lfloor n/d \rfloor}{\lfloor k/d \rfloor} \cdot
\qbinom{n \bmod d}{k \bmod d}\!\bigg|_{t = \zeta},
\]
where the first factor is an \emph{ordinary} binomial coefficient and the
second is a $q$-binomial \emph{evaluated at $\zeta$}.
\end{lemma}

Specialising to $d = 3$ and $\zeta = \omega$, the possible inner-factor values are
\[
\qbinom{0}{0}\!\!\bigg|_\omega \!\!= 1,\quad
\qbinom{1}{0}\!\!\bigg|_\omega \!\!= \qbinom{1}{1}\!\!\bigg|_\omega \!\!= 1,\quad
\qbinom{2}{0}\!\!\bigg|_\omega \!\!= \qbinom{2}{2}\!\!\bigg|_\omega \!\!= 1,\quad
\qbinom{2}{1}\!\!\bigg|_\omega \!\!= 1 + \omega = -\omega^2,
\]
and $\qbinom{r}{s}\!|_\omega = 0$ whenever $s > r$ within $\{0, 1, 2\}$.

Denote $B_r := (-1)^r \qbinom{2 \bmod 3}{r}|_\omega$, the sign-weighted
inner-factor values used when $2k \equiv n \bmod 3$. There are three
regimes depending on $n \bmod 3 \in \{0, 1, 2\}$:

\begin{itemize}
\item ($n \equiv 0$): $B_0 = 1$, $B_1 = B_2 = 0$.
\item ($n \equiv 1$): $B_0 = 1$, $B_1 = -1$, $B_2 = 0$. \textbf{The full-triple sum $B_0 + B_1 + B_2 = 0$.}
\item ($n \equiv 2$): $B_0 = 1$, $B_1 = -(-\omega^2) = \omega^2$, $B_2 = 1$; full-triple sum $= 2 + \omega^2 = 1 - \omega \ne 0$.
\end{itemize}

Since $n = 2k$: $k \equiv 0 \pmod 3 \Leftrightarrow n \equiv 0$;
$k \equiv 1 \pmod 3 \Leftrightarrow n \equiv 2$;
$k \equiv 2 \pmod 3 \Leftrightarrow n \equiv 1$.

\section{Proof of Theorem \ref{thm:main}}

By Proposition~\ref{prop:telescope}, both $m_{(2)}(t)$ and
$m_{(1,1)}(t)$ are, up to sign, of the form $S_j(t)$ with $j \in
\{k-1, k\}$.  It suffices to show
\begin{enumerate}
\item[(a)] $\qbinom{2k}{k}|_\omega = 0 \iff k \equiv 2 \pmod 3$;
\item[(b)] at $k \equiv 2 \pmod 3$, $m_{(2)}(\omega) = 0$;
\item[(c)] at $k \equiv 0 \pmod 3$ or $k \equiv 1 \pmod 3$, $m_{(2)}(\omega) \ne 0$.
\end{enumerate}

Combining (a) with (b) gives $m_{(2)}(\omega) = m_{(1,1)}(\omega) = 0$ (using the sum identity), so both $[3]_t$-divide; that's the $\Leftarrow$ direction of Theorem~\ref{thm:main}. Item (c) gives the $\Rightarrow$ direction (at least $m_{(2)}$ fails).

\medskip
\emph{Proof of (a).} By Lemma~\ref{lem:qlucas},
$\binom{2k}{k}_\omega = \binom{\lfloor 2k/3 \rfloor}{\lfloor k/3 \rfloor} \cdot \qbinom{2k \bmod 3}{k \bmod 3}\!|_\omega$.
The inner factor is $\qbinom{0}{0}|_\omega = 1$ ($k \equiv 0$),
$\qbinom{2}{1}|_\omega = -\omega^2 \ne 0$ ($k \equiv 1$), or
$\qbinom{1}{2}|_\omega = 0$ ($k \equiv 2$). Hence the product is zero
iff $k \equiv 2 \pmod 3$.

\medskip
\emph{Proof of (b).} Write $k = 3s + 2$, $n = 2k = 6s + 4 \equiv 1 \pmod 3$; so
$\lfloor n/3 \rfloor = 2s + 1$ and $n \bmod 3 = 1$.
By Lemma~\ref{lem:qlucas},
\[
c_i|_\omega = \qbinom{2k}{i}\!\bigg|_\omega
= \binom{2s+1}{\lfloor i/3 \rfloor} \cdot \qbinom{1}{i \bmod 3}\!\bigg|_\omega.
\]
The inner factor is $1$ if $i \bmod 3 \in \{0, 1\}$ and $0$ if $i \bmod 3 = 2$.
Splitting $i = 3q + r$,
\[
S_j(\omega) \;=\; \sum_{q : 3q \le j} (-1)^q \binom{2s+1}{q} \cdot
\underbrace{\bigl[ B_0 \cdot \mathbb{1}[3q \le j] + B_1 \cdot \mathbb{1}[3q+1 \le j] \bigr]}_{=: E_q},
\]
where $B_0 = 1$, $B_1 = -1$ (the $B_2$-term is absent since the inner factor $\qbinom{1}{2}|_\omega = 0$).

For any $q$ with $3q + 1 \le j$, $E_q = B_0 + B_1 = 0$; only the ``last'' $q$ (where the constraint tightens) can contribute.

For $\rho := j - 3\lfloor j/3\rfloor \in \{0, 1, 2\}$:
\begin{itemize}
\item $\rho = 0$: at $q = j/3$, $3q = j$ but $3q + 1 = j + 1 > j$; $E_q = B_0 = 1$; contribution $(-1)^{j/3} \binom{2s+1}{j/3}$.
\item $\rho = 1$: at $q = (j-1)/3$, both $3q = j-1 \le j$ and $3q+1 = j \le j$; $E_q = B_0 + B_1 = 0$. \emph{No contribution.}
\item $\rho = 2$: at $q = (j-2)/3$, both $3q + 1 \le j$; $E_q = 0$. \emph{No contribution.}
\end{itemize}

Now split by parity of $k = 3s+2$:
\begin{itemize}
\item $k$ even $\Longrightarrow s$ even, so $k = 6t + 2$; $j = k = 6t+2$; $\rho = 6t+2 \bmod 3 = 2$. Contribution zero: $S_j(\omega) = 0$.
\item $k$ odd $\Longrightarrow s$ odd, so $k = 6t + 5$; $j = k - 1 = 6t + 4$; $\rho = (6t+4) \bmod 3 = 1$. Contribution zero: $S_j(\omega) = 0$.
\end{itemize}

Either way, $m_{(2)}(\omega) = S_j(\omega) = 0$. \qed (b)

\medskip
\emph{Proof of (c).} Two subcases.

\smallskip
\emph{Case $k \equiv 0 \pmod 3$.} Write $k = 3s$, $n = 2k = 6s \equiv 0 \pmod 3$; $\lfloor n/3 \rfloor = 2s$.
Inner factor $\qbinom{0}{r}|_\omega$ is $1$ if $r = 0$ else $0$. Hence
$c_i|_\omega = \binom{2s}{i/3}$ if $3 \mid i$, and $0$ otherwise. Thus
\[
S_j(\omega) \;=\; \sum_{q=0}^{\lfloor j/3 \rfloor} (-1)^{3q} \binom{2s}{q}
\;=\; \sum_{q=0}^{Q} (-1)^{q} \binom{2s}{q},
\]
with $Q = \lfloor j/3\rfloor$. The classical alternating identity
$\sum_{q=0}^{Q}(-1)^q \binom{N}{q} = (-1)^Q \binom{N-1}{Q}$ (for $Q < N$)
gives:
\begin{itemize}
\item $k$ even $= 6u$: $j = k = 6u$; $Q = 2u$; $2s = 4u$; $S_j(\omega) = \binom{4u-1}{2u}$. Positive for $u \ge 1$.
\item $k$ odd $= 6u + 3$: $j = k-1 = 6u+2$; $Q = 2u$; $2s = 4u+2$; $S_j(\omega) = \binom{4u+1}{2u}$. Positive for $u \ge 0$.
\end{itemize}
Both are positive integers, so nonzero, so $m_{(2)}(\omega) \ne 0$.

\smallskip
\emph{Case $k \equiv 1 \pmod 3$.} Write $k = 3s + 1$, $n = 2k = 6s+2 \equiv 2 \pmod 3$; $\lfloor n/3 \rfloor = 2s$.
Inner factor: $B_0 = 1$, $B_1 = \omega^2$ ($= -(-\omega^2)$),
$B_2 = 1$. Full-triple sum: $B_0 + B_1 + B_2 = 2 + \omega^2 = 1 - \omega$
(using $1 + \omega + \omega^2 = 0$).

Split $i = 3q + r$: $c_i|_\omega = \binom{2s}{q} \cdot \qbinom{2}{r}|_\omega$.
For $q$ with $3q + 2 \le j$ (all three $r$'s in range): $E_q^{\text{full}} = 1 - \omega$.
For the boundary $q$:
\begin{itemize}
\item $\rho = 0$: $E_Q = 1$;
\item $\rho = 1$: $E_Q = 1 + \omega^2 = -\omega$;
\item $\rho = 2$: full ($E_Q = 1 - \omega$).
\end{itemize}

Split by parity of $k = 3s+1$:
\begin{itemize}
\item $k$ even $= 6u + 4$ (so $s = 2u+1$, $2s = 4u+2$): $j = k = 6u+4$; $Q = 2u+1$; $\rho = 1$.
\begin{align*}
S_j(\omega) &= (1 - \omega)\sum_{q=0}^{2u} (-1)^q \binom{4u+2}{q} + (-1)^{2u+1}\binom{4u+2}{2u+1} \cdot (-\omega) \\
&= (1 - \omega)\binom{4u+1}{2u} + \omega \binom{4u+2}{2u+1}.
\end{align*}
Now $\binom{4u+2}{2u+1} = 2\binom{4u+1}{2u}$ (Pascal's rule $\binom{N+1}{K+1} = \binom{N}{K} + \binom{N}{K+1}$ with $N = 4u+1, K = 2u$ and symmetry $\binom{4u+1}{2u} = \binom{4u+1}{2u+1}$), so
\[
S_j(\omega) = \binom{4u+1}{2u}\bigl[(1-\omega) + 2\omega\bigr] = \binom{4u+1}{2u}(1 + \omega) = -\omega^2 \binom{4u+1}{2u}.
\]
For $u \ge 0$: $\binom{4u+1}{2u} \ge 1$, so $|S_j(\omega)| = \binom{4u+1}{2u} > 0$.

\item $k$ odd $= 6u + 1$ (so $s = 2u$, $2s = 4u$): $j = k-1 = 6u$; $Q = 2u$; $\rho = 0$.
\begin{align*}
S_j(\omega) &= (1 - \omega)\sum_{q=0}^{2u-1} (-1)^q \binom{4u}{q} + (-1)^{2u}\binom{4u}{2u} \cdot 1 \\
&= -(1-\omega)\binom{4u-1}{2u-1} + \binom{4u}{2u}.
\end{align*}
Using $\binom{4u}{2u} = 2\binom{4u-1}{2u-1}$ (Pascal / symmetry as above),
\[
S_j(\omega) = \binom{4u-1}{2u-1}\bigl[-(1-\omega) + 2\bigr] = \binom{4u-1}{2u-1}(1+\omega) = -\omega^2 \binom{4u-1}{2u-1}.
\]
For $u \ge 1$: nonzero. For $u = 0$ (i.e., $k = 1$): $S_0(\omega) = c_0 = 1 \ne 0$ directly.
\end{itemize}
So $S_j(\omega) = -\omega^2 \cdot (\text{positive integer}) \ne 0$ in all cases.

\medskip
In both subcases (c), $m_{(2)}(\omega) \ne 0$; hence $[3]_t \nmid m_{(2)}(t)$.
Theorem~\ref{thm:main} is proved. \qed

\section{Computational verification}

The following are verified by SymPy in
\texttt{\~{}/projects/probes/2026-07-31-r2-proof-verification/verify.py}:

\begin{enumerate}
\item Reduction identity of Proposition~\ref{prop:telescope}: $m_{(2)}(t)$
and $m_{(1,1)}(t)$ computed directly via
Theorem~\ref{thm:D}'s isotypic sum vs.\ via $\pm S_j(t)$ agree as polynomials
in $\mathbb{Z}[t]$ for all $k \in \{1, \ldots, 8\}$ (zero mismatches).
\item Divisibility pattern: for $k \in \{1, \ldots, 12\}$,
$[3]_t \mid m_\pi(t)$ for both $\pi \in \{(2), (1,1)\}$ iff
$k \equiv 2 \pmod 3$ (verified by polynomial long division; zero mismatches).
\item The $q$-Lucas evaluation of $\qbinom{2k}{i}|_\omega$ agrees
symbolically with Lemma~\ref{lem:qlucas} for $k \in \{2, \ldots, 6\}$
and all $i \in \{0, \ldots, 2k\}$ (all direct computations $=$ the
ordinary-binomial-times-inner-factor product).
\item The closed-form values of $S_j(\omega)$ from the proof of (c)
agree numerically with the direct SymPy evaluation
for $k \in \{1, \ldots, 12\}$, up to the phase factor $-\omega^2$ (which
has $|-\omega^2| = 1$).
\end{enumerate}

\section{Discussion and outlook}

The proof structure is:
\begin{center}
\emph{isotypic sums} $\xrightarrow[\text{Lemma 2}]{\text{two-row fake deg}}$
\emph{telescope of $q$-binomials} $\xrightarrow[\text{Lemma 4}]{q\text{-Lucas}}$
\emph{finite $\omega$-arithmetic}.
\end{center}
The vanishing in case (b) is not the vanishing of any individual
$\fdeg_\lambda(\omega)$ (indeed we know from the probe that the
individual fake degrees typically do \emph{not} divide by $[3]_t$); it
is a coherent cancellation induced by the alternating structure of the
telescope and by the fact that the ``inner'' full-triple
$q$-binomial sum $\qbinom{1}{0}|_\omega - \qbinom{1}{1}|_\omega + 0 = 0$
vanishes precisely when $n \equiv 1 \pmod 3$ (i.e., $k \equiv 2$).

The pattern of the survivor in case (c) is striking: the residual value
is always $-\omega^2 \cdot \binom{2s' \pm 1}{s'}$ for a central binomial
coefficient. Both cases $k \equiv 0$ and $k \equiv 1$ produce
central-binomial magnitudes with real vs.\ complex-$\zeta_3$-phase, and
the magnitudes are all positive integers.

\begin{remark}[Generalisation to $r \ge 3$]
The full conjecture (\cite{PROVE-2026-07-31}) states
$[2r-1]_t \mid m_\pi(t)$ for every $\pi \vdash r$ iff
$k \bmod (2r-1) \in \{2, \ldots, 2r-2\}$. The $r = 2$ case above is the
smallest instance. For $r \ge 3$ the telescoping argument does not
generalise directly; one likely route is Springer's theorem on regular
elements applied to the wreath $S_k \wr S_r$, or the wreath-plethystic
formula of Macdonald Ch~I App~B combined with the fake-degree formula at
$\zeta_d$ (Désarménien).
\end{remark}

\bibliographystyle{alpha}
\begin{thebibliography}{9}
\bibitem{Wreath-r=2} C.\ Muse, \emph{Wreath decomposition of the collision
regime coset Poincar\'e}, note (2026-07-30).
\bibitem{Wreath-general-r} C.\ Muse, \emph{The $r$-wreath extension of
Theorem D}, note (2026-07-30).
\bibitem{PROVE-2026-07-31} \texttt{\~{}/state/PROVE.md.completed-20260731},
period-$(2r-1)$ divisibility conjecture.
\end{thebibliography}

\end{document}
